% Main-text fragment; requires graphicx. Compile from docs/paper/.
% Include online_appendix.tex after the paper's appendix declaration.
% Do not also include the online figure block from task_illustrations.tex.
\subsection{Continuous online reasoning and execution}
\label{sec:online-experiments}

We next test whether geometric abstraction helps an LLM use prior skills and
complete tasks under a real-time budget. The agent repeatedly observes, reasons,
and acts in one continuing scene (Figure~\ref{fig:online-task-illustrations}).
An unsuccessful action leaves its consequences in place; the agent must choose
its next action from the resulting state. Model waiting time therefore consumes
the same budget as physical execution and recovery.

\paragraph{Tasks and shared skills.}
We use three SO-100 tasks. \emph{Bridge-tower} composes bridge construction with
stacking two cubes above the beam. \emph{Ring-on-peg} requires aligning a grasped
ring with a peg and releasing it. \emph{Tower relocation} transfers the upper
two cubes of an existing column onto another foundation, preserving their order
and the foundation positions. Both conditions receive the same prior skills:
adapted bridge/stacking routines, a trained ring grasp-and-lift prefix, and
single-placement routines for relocation. The Ring prefix is selected using
training outcomes and uses local geometry in both conditions; it stops before
placement. The online agent chooses targets, action order, additional observations,
and recovery. Skills execute locally without internal LLM calls.

\paragraph{Comparison and episode protocol.}
Base and Full use separate persistent \texttt{gpt-6-astra} contexts with medium
reasoning effort. They share task goals, prior skills, and action interfaces.
Both can request images and calibrated pixel-depth samples. Base receives images
and qualitative feedback; Full additionally receives compact current geometric
text and references, with initial and requested images available. Shared actions
execute the submitted targets. Runtime inputs exclude simulator object truth and
private evaluator state. Appendix~\ref{app:geometric-abstraction} details the
visual grounding and text interface.

Each task uses 50 paired layouts, with one selected episode per condition/layout
(300 episodes in the analysis). This combines the original five-pair pilot with a
separately registered extension of 45 new pairs per task. The environment resets
once at initialization. For the
next 200 seconds, the agent alternates decisions and execution until verified
success or timeout, without resetting after failures. Physics continues during
model waits. The clock includes initialization after reset, local sensing,
model requests, skill execution, and settling; the Ring grasp/lift is included.
There are no injected disturbances. Success requires the target arrangement,
release, and a stable hold. We report success counts and mean capped completion
time, assigning 200 seconds to timeouts. Model time and recorded tokens are
cumulative per-episode costs averaged over all 50 episodes, including failures.
Both arms of nine layouts affected by timing/GPU faults were repeated once;
all audited repeats replace those original attempts. The new cohort, original
outcomes, and a sensitivity excluding these pairs are reported in
Appendix~\ref{app:online-details}.

\begin{figure*}[t]
  \centering
  \includegraphics[width=\textwidth]{figures/online_workflow_tasks.png}
  \caption{\textbf{Online execution and tasks.} The agent observes and acts in one
  scene until success or timeout. Full receives geometric abstraction for current
  decisions. Right: illustrative starts and goals; Ring is shown after the shared
  lift. Task panels are constructed renderings, not recorded trajectories.}
  \label{fig:online-task-illustrations}
\end{figure*}

\begin{figure*}[t]
  \centering
  \includegraphics[width=\textwidth]{figures/outcome_bars.png}
  \caption{\textbf{Online outcomes and cost.} Success within 200 seconds (left),
  mean capped completion time (middle), and mean cumulative recorded tokens per
  episode (right). Each task uses 50 paired layouts. All episodes enter
  the means; timeouts count as 200 seconds. Daggers mark unknown additional
  interrupted token usage, excluded from recorded totals. Nine pairs use audited
  repeats; paired intervals and sensitivity appear in the appendix.}
  \label{fig:online-outcomes}
\end{figure*}

\paragraph{More completions and earlier success.}
Full completes more episodes on Bridge-tower (28/50 versus 12/50) and Relocation
(42/50 versus 15/50), with lower mean capped time on both
(Figure~\ref{fig:online-outcomes}). The paired success differences are 32 and
54 percentage points, with 95\% paired-bootstrap intervals of [16, 48] and
[34, 72], respectively. Ring approaches ceiling success in both conditions
(50/50 Full, 49/50 Base); its main benefit is earlier completion. Full has lower
capped time on 45/50 pairs, and mean capped time falls from 115.8 to 77.2 seconds,
a 33.3\% reduction. Relocation's mean capped time falls from 186.4 to 139.2
seconds. The added layouts support the pilot's pattern of more construction
completions and faster Ring placement under the same wall-clock budget.

\paragraph{Less repeated model interaction.}
Figure~\ref{fig:online-model-resources} shows cumulative model-request time
and recorded tokens, together with variation across episodes. On Ring, Full
reduces cumulative model time by 49.4\% and recorded tokens by
52.8\%. Completed requests fall from 11.88 to 6.90 per episode, and image
attachments from 10.94 to 2.72 per episode
(Table~\ref{tab:online-results}, Appendix~\ref{app:online-details}).
Relocation shows a similar pattern: model time
falls by 60.5\% and tokens by 52.5\%, while accumulated physical-action time is
similar (103.2 versus 104.5 seconds). Bridge-tower's higher success instead comes
with 15.1\% more recorded tokens and only 6.7\% lower cumulative model time.
Reduced token consumption is therefore task-dependent.

\begin{figure*}[t]
  \centering
  \includegraphics[width=\textwidth]{figures/model_resources.png}
  \caption{\textbf{Model interaction cost.} Bars are means over all 50 episodes,
  including timeouts; points show individual episodes. Left: cumulative request
  wall time, including interrupted waits. Right: recorded input-plus-output tokens,
  counting cached input once. Triangles identify unknown additional interrupted
  usage, excluded from the plotted values.}
  \label{fig:online-model-resources}
\end{figure*}

Full acquires more local camera frames but attaches fewer images to the LLM on
all three tasks. This is consistent with the intended role of geometric
abstraction: local perception supplies actionable state, reducing repeated
visual interpretation during interaction. The benefit concerns the complete
reasoning-and-execution loop, including action choices and request frequency.
The expansion measures variation across more layouts, with one selected episode
per condition/layout. It still compares the supplied geometry and feedback package
with matched skills; it does not isolate measurement quality, symbolic references,
text format, or model inference speed. Appendix~\ref{app:online-details} gives
cohort results, paired intervals, runtime sensitivities, and success-versus-time
curves.
